% =====================================================================
% Slide 21: failure analysis & path forward
% =====================================================================
\begin{frame}{Failure Analysis \& Path Forward}
\begin{columns}[T]
\column{0.49\textwidth}
\textbf{Why does Stage~5b fail today?}
\begin{itemize}\small
  \item \textbf{Identifiability.} The cross-term block of a
        12-element MEAM exposes $\sim\!200$ free parameters
        against $\le\!100$ medoid-selected
        composition targets --- the inverse problem is
        under-determined.
  \item \textbf{Surrogate--LAMMPS mismatch.} The NN approximates
        the LAMMPS-elastic map with residual RMSE
        $\sim\!10\,$GPa; the L-BFGS optimizer chases artefacts of
        this residual into MEAM regions where the true LAMMPS
        gradient bends sharply.
  \item \textbf{Bound-saturation.} Several optimized parameters
        hit their box bounds (esp.\ $E_c(i,j)$, $\alpha(i,j)$),
        which is a tell-tale of an objective that wants to leave
        the physically meaningful region.
\end{itemize}

\column{0.49\textwidth}
\textbf{Planned next steps}
\begin{itemize}\small
  \item \textbf{Parameter-subset selection.} Freeze the
        single-element blocks (initialised by Stage~4),
        refit only the most-sensitive cross-terms identified by
        Sobol sensitivity.
  \item \textbf{Active learning expansion.}
        \texttt{--sampling active} already implemented;
        scale to 500+ medoids using ensemble-disagreement
        acquisition.
  \item \textbf{Physics-aware loss.} Add penalties for
        $\nu \!\to\! 0.5$ and $C_{11}\!<\!C_{12}$
        \emph{inside} the surrogate, not just at the
        verification gate.
  \item \textbf{Tighter QE coverage.} Add intermediate-composition
        SQS~\citep{lewis1995} DFT points so the cross-term loss can be
        evaluated against \emph{multiple} reference compositions per
        pair.
\end{itemize}
\end{columns}
\end{frame}

% =====================================================================
% Slide 22: Summary
% =====================================================================
\begin{frame}{Summary}
\begin{block}{Delivered}
\scriptsize
\begin{itemize}\itemsep0.02em
  \item \textbf{End-to-end pipeline} (Stages 1--4 + 5a):
        random alloy $\to$ MD elastic $\to$ DFT reference $\to$
        merged 12-element MEAM library $\to$ composition surrogate.
  \item \textbf{Composition NN surrogate} with deep-ensemble
        uncertainty, $R^2 \!\sim\! 0.65$ on held-out $C_{11}$,
        $R^2\!\sim\!0.77$ on derived $E$, distributed as an
        interactive predictor.
  \item \textbf{Closed-loop NN-MEAM optimizer} (Stage 5b) with
        atomic checkpointing, three sampling regimes (Sobol, Latin
        hypercube, and active learning) and an honest verification
        protocol.
\end{itemize}
\end{block}

\begin{block}{Methodological contribution}
\scriptsize
A closed-loop framework in which a \emph{neural-network surrogate
of the MEAM-to-stiffness map} is inverse-optimized against
MD-measured targets --- enabling MEAM construction from
first-principles input alone, with no manual fitting and no
experimental data. Verification of the NN-driven Stage~5b on the
current 12-element corpus is the open problem to attack next.
\end{block}
\end{frame}
